\section{A rational-base surface correspondence with zero transcendental action}
\label{sec:mumford-fiber-product}

This section tests a bounded surface correspondence coming directly
from the trigonal equation of a genus-four Mumford family.
It produces a smooth K3 surface and a correspondence of degrees
two and nine, but its transcendental action vanishes.
The general reason is factorization through a rational surface.

\begin{lemma}\label{lem:rational-base-fiber-product}
Let $S$ and $X$ be smooth projective complex surfaces, let $R$ be
a smooth projective rational surface, and let $p:S\dashrightarrow R$
be dominant.  Suppose $f:X\to R$ is finite and flat.
Resolve $p$ to $q:S'\to R$.  The full fibre-product cycle
$S'\times_R X$, pushed to $S\times X$, acts as zero on
$T(S)=\operatorname{NS}(S)_{\mathbf Q}^{\perp}$.
If $p$ has degree two and $f$ has odd degree, that fibre product
has just one dominant irreducible component and no additional
surface components over a proper subset of $S'$.
\end{lemma}

\begin{proof}
Flatness gives the expected graph-composition cycle, whose action is
$f^*q_*:H^2(S',\mathbf Q)\to H^2(X,\mathbf Q)$.
Every class of $H^2(R,\mathbf Q)$ is a divisor class.
For $\alpha\in T(S')$ and $\eta\in H^2(R,\mathbf Q)$,
\[
 \int_R(q_*\alpha)\eta=\int_{S'}\alpha q^*\eta=0.
\]
Nondegeneracy of the intersection form gives $q_*\alpha=0$.
Point blowups add only divisor classes, and therefore identify
$T(S)$ with $T(S')$.  This proves the first assertion.

A quadratic function-field extension remains a field after an
odd-degree extension: otherwise the latter would contain a
quadratic subextension, contrary to the tower formula.
This proves generic irreducibility.  The base change to $S'$ is
finite flat, so every top-dimensional component dominates $S'$;
flatness excludes components with vertical support.  This proves
the remaining assertions.
\end{proof}

The lemma concerns the full fibre-product cycle.  If an even-degree
base change splits a double cover, an individual component could
behave differently; finding that component requires a rational
lift to $S$.  The odd-degree construction below offers no such
choice.

\subsection{An explicit trigonal curve and elliptic K3 surface}

Take the family in
\cite[Theorem~1.2(ii), equation~(1.2)]{BHPS}.
On the quadric cone $XZ=Y^2$, use the chart
$Z=1$, $X=x^2$, $Y=x$, $W=w$.  Its equation becomes
\[
 C_t:\quad w^3+A_t(x)w+B_t(x)=0,
\]
where
\begin{align*}
 a_t(x)&=5x^4+6x^3+2tx+3t,\\
 b_t(x)&=(9-2t)x^6+22tx^5+21tx^4+(-14t^2+18t)x^3\\
       &\hspace{2em}+t^2x^2+6t^2x+3t^2(2-t),\\
 A_t(x)&=t(t-1)a_t(x),\qquad
 B_t(x)=\frac{t(t-1)}3b_t(x).
\end{align*}
Put $\Delta_t=-4A_t^3-27B_t^2$.  Direct expansion gives
\begin{equation}
 \Delta_t(0)=-27t^8(t-1)^2,
 \qquad [x^{12}]\Delta_t=-t^2(t-1)^2(512t^2-608t+243).
 \label{eq:bhps-discriminant-endpoints}
\end{equation}
There is a nonempty open set of parameters for which $t\ne0,1$,
both quantities are nonzero, and $\Delta_t$ has twelve simple roots.
Here is an exact certificate of nonemptiness.  At $t=2$ write
$\Delta_2=4d$.  Reduction modulo seven gives
\begin{align*}
 d={}&3x^{12}+x^{11}+x^{10}+2x^7+4x^6+4x^5+2x+1
                    \quad\text{in }\mathbf F_7[x].
\end{align*}
For
\begin{align*}
 u&=-2x^{10}-x^9-x^8-2x^7-2x^6-2x^2+3,\\
 v&=-x^{11}+2x^{10}+x^9-2x^8+x^5+3x^4+2x^3+x^2-3x-1,
\end{align*}
the identity $ud+vd'=1$ holds in $\mathbf F_7[x]$.
The leading coefficient is nonzero, so $\Delta_2$ is squarefree
of degree twelve in characteristic zero.  The
\emph{trigonal fibre-product check} in the
\href{\workbenchurl}{workbench} verifies this polynomial identity
and \eqref{eq:bhps-discriminant-endpoints} by exact arithmetic.

For a parameter in this open set, the projective curve $C_t$ is
smooth.  A simple discriminant zero gives one simple ramification
point and no singularity; a triple root would instead give a
multiple discriminant zero.  At infinity set $z=1/x$ and $w'=z^2w$.
The equation becomes
\[
 w'^3+z^4A_t(1/z)w'+z^6B_t(1/z)=0,
\]
whose discriminant is a unit at $z=0$ by
\eqref{eq:bhps-discriminant-endpoints}.
The vertex of the quadric cone is not on the cubic.
These charts thus cover the curve and show that the ruling
$x:C_t\to\mathbf P^1$ is unramified over zero and infinity.
The Koszul resolution of the complete intersection of degrees two
and three in $\mathbf P^3$ gives $H^0(\mathcal O_{C_t})=\mathbf C$.
It is therefore connected and, being smooth, irreducible.
Adjunction gives canonical degree six, so it has genus four.
The ruling is finite flat of degree three.

Consider the smooth relatively minimal elliptic surface with section
defined by the Weierstrass equation
\begin{equation}
 S_t:\quad Y^2=X^3+x^2A_t(x)X+x^3B_t(x).
 \label{eq:bhps-elliptic-k3}
\end{equation}
The coefficients are sections of $\mathcal O(8)$ and $\mathcal O(12)$,
so its fundamental line bundle is $\mathcal O(2)$.
At the twelve roots of $\Delta_t$ the elliptic discriminant has
simple zeros, giving twelve fibres of type $I_1$.
At zero the coefficient valuations are at least two and three,
the first is exactly two, and the discriminant valuation is six.
This is a minimal $I_0^*$ fibre.
At infinity use $X'=z^4X$, $Y'=z^6Y$.
The transformed coefficients are $z^6A_t(1/z)$ and $z^9B_t(1/z)$,
again of valuations two and at least three, with discriminant
valuation six.  This gives a second minimal $I_0^*$ fibre.
The minimality and fibre criteria are those of
\cite[\S4.5, Table~1]{SchuttShioda}.
Thus the fibre configuration is $2I_0^*+12I_1$.

The canonical bundle formula
\cite[Theorem~6.8]{SchuttShioda} gives
$\omega_{S_t}=\pi^*(\mathcal O(-2)\otimes\mathcal O(2))
 =\mathcal O_{S_t}$.
Relative duality gives $R^1\pi_*\mathcal O_{S_t}=\mathcal O(-2)$,
and Leray gives $H^1(S_t,\mathcal O)=0$.
Thus $S_t$ is K3.  Its fibre Euler numbers add to $24$, also
consistent with \cite[Corollary~6.11]{SchuttShioda}.
The section, fibre and the two root configurations give
\[
 U\oplus D_4(-1)\oplus D_4(-1)
                       \subset\operatorname{NS}(S_t).
\]
This shows Picard rank at least ten.  No identification with the
Picard-thirteen real-multiplication surface of
Theorem~\ref{thm:mumford-surface} is asserted.

\subsection{The irreducible correspondence of degrees two and nine}

On $x\ne0$, put $w=X/x$ and $y=Y/x$ in
\eqref{eq:bhps-elliptic-k3}.  Then
\[
 y^2=x\bigl(w^3+A_t(x)w+B_t(x)\bigr).
\]
The coordinates $(x,w)$ give a rational double cover
$S_t\dashrightarrow R=\mathbf P^1_x\times\mathbf P^1_w$.
Resolve it to $q:S'_t\to R$.
Set $X_t=C_t\times C_t$ and define
\[
 f:X_t\longrightarrow R,\qquad
       (p_1,p_2)\longmapsto(x(p_1),x(p_2)).
\]
Here the second trigonal coordinate is used as the $w$ coordinate
of $R$.  This is finite flat of degree nine.
The fibre product has, over $X_t$, the generic equation
\[
 y^2=x_1\bigl(x_2^3+A_t(x_1)x_2+B_t(x_1)\bigr).
\]
At a divisor $\{P\}\times C_t$ with $x(P)=0$, the first factor
has valuation one because the ruling is unramified there.
The second factor restricts to a nonzero polynomial in the
nonconstant function $x_2$, so the radicand has odd valuation.
It is not a square.  This also checks generic irreducibility
directly, in addition to the odd-degree argument of the lemma.

Finite flatness over $S'_t$ excludes vertical surface components.
Normalize and resolve the unique surface component, then push its
fundamental class to $S_t\times X_t$.  The degrees of its two
projections are respectively nine and two.
All equations have fixed degrees on the open parameter set,
so these correspondences form a bounded algebraic family.
Nevertheless Lemma~\ref{lem:rational-base-fiber-product} gives
\[
 [Z_t]_*|_{T(S_t)}=0.
\]
In particular its component from $T(S_t)$ to
$H^1(C_t)\otimes H^1(C_t)$ is zero.
The transpose also has zero projection to $T(S_t)$ by adjointness.
Normalization, resolution and the blowups of $S_t$ do not change
this conclusion: their exceptional contributions are divisor
classes, and the fundamental fibre-product correspondence still
acts as $f^*q_*$.
This rules out the indicated correspondence even if a later
identification of the K3 surface's transcendental structure were
available.  It leaves arbitrary surface correspondences and the
coniveau condition of Theorem~\ref{thm:mumford-coniveau} open.
